%=========================================================================
% Template authors: Michal Bidlo, Bohuslav Křena, Jaroslav Dytrych, Petr Veigend and Adam Herout 2019
% Thesis author: Matúš Csirik 2025

\chapter{Introduction}\label{ch:intro}
% Motivation, problem context, thesis goals, scope, contributions overview, structure summary.
% list explicit objectives

\chapter{Background on Automated Assessment of Technical Documentation}\label{ch:background}
% Mention interesting problems faced in the landscape.

\section{Figures, Text, and Layout Extraction}\label{sec:extraction}
% work, tools, PDF/MD parsing, vision models, layout analysis, limitations.
\subsection*{PDF Document Parsing}\label{subsec:doc-parsing-pdf}
\subsection*{Markdown Document Parsing}\label{subsec:doc-parsing-markdown}
\subsection*{Source-Code Repository Analysis}\label{subsec:source-code-analysis}
\subsection*{Figure Analysis}\label{subsec:figure-analysis}

\section{Human-in-the-Loop Grading}\label{sec:hitl}

\section{Risks and Design Takeaways}\label{sec:risks}
% bias, reliability, privacy, requirements


\chapter{Current State and Needs in Local Courses}\label{ch:local-state}

\section{Workflows and Roles}\label{sec:workflows}
% students, graders, supervisors, current process mapping

\section{Assessment-Criteria Codebook}\label{sec:codebook}
% Criterion list, definitions, rationale, ambiguity 

\section{Gap Analysis}\label{sec:gap-analysis}
% Pain points vs capabilities, show opportunities for improvements.

\section{Objectives and Scope}\label{sec:objectives}
% Concrete objectives, in/out of scope decisions.


\chapter{System Design and Prototype Development}\label{ch:design}
% Architectural goals, relate to requirements from risks and gap-analysis.

\section{Data Flow Overview}\label{sec:dataflow}
% End-to-end pipeline narrative (input -> parse -> detect -> aggregate -> report).

\section{Document and Source-Code Representation}\label{sec:representation}
% Internal IR JSON, pydantic, normalization rules, source->IR mapping.

\section{Detectors for Assessment-Criteria Coverage}\label{sec:detectors}
% Overview table: detector id, AC codes, feature sources, output fields.
\subsection*{detector1}\label{subsec:detector1}
% Purpose, inputs, algorithm/heuristic/model, output schema, limitations, iteration...
\subsection*{detector2}\label{subsec:detector2}
\subsection*{detectorx}\label{subsec:detectorx}
% (placeholder list).

\section{Rule Engine and Decision Logic}\label{sec:rule-engine}
% Rule representation, conflict resolution, threshold handling.

\section{Evidence, Reporting, and Traceability}\label{sec:evidence}
% linking to IR spans, reproducibility

\section{Configuration Presets}\label{sec:presets}
% IFJ vs IPP parameter differences.

\section{Major Design Iterations and Rationale}\label{sec:iterations}
% Timeline of rejected alternatives and why. if relevant


\chapter{Model Evaluation}\label{ch:evaluation}
% Evaluation philosophy of the entire ensemble model.

\section{Success Criteria}\label{sec:success-criteria}
% Quantitative + qualitative targets (accuracy, coverage, reviewer time saved, etc.).

\section{Provided Datasets and Experimental Setup}\label{sec:datasets}
% Dataset description, reproducibility steps.

\section{Aggregate Results and Analysis}\label{sec:results}
% Core metrics tables, analysis commentary.

\section{Agreement and Baseline Comparison}\label{sec:agreement}
% baseline heuristic comparison with human grading, statistical significance.

\section{Cost, Latency, and Throughput}\label{sec:cost}
% Performance (per doc, per detector), scaling considerations.

\section{Error Analysis and Generalization}\label{sec:error}
% Errors, generalization risks, mitigations.


\chapter{Pilot Deployment}\label{ch:pilot}

\section{Deployment Plan and Environment}\label{sec:deployment}
% Environment, integration steps, timeline, risk mitigations.
\section{Integration Workflow and Configuration Presets}\label{sec:integration}
% Operational flow, preset selection and overrides.
\section{Pilot Results}\label{sec:pilot-results}
% Observed metrics vs expectations, feedback.
\section{Fairness and Operational Trade-Offs}\label{sec:fairness}
% Bias evaluation, false positive/negative impact by criterion.
\section{Lessons Learned and Next Steps}\label{sec:lessons}
% Planned enhancements and future research directions.

\chapter{Conclusion}\label{ch:conclusion}
% Recap objectives, achieved contributions, limitations, future work roadmap.
%=========================================================================